\documentclass[review,3p,times]{elsarticle}
\usepackage{url}
\begin{document}

\section*{Research in Context}

\subsection*{Systematic review}
Motor dysfunction, impaired manual dexterity, and motor-cortex pathology have been reported in Alzheimer's disease (AD). However, many MRI analyses of early motor-system involvement rely on cognitive status or cognitive scores rather than PET biomarkers to define AD pathophysiology.

\subsection*{Interpretation}
Using OASIS-3 amyloid PET, adults with MMSE $>24$ showed exploratory cortical-thickness associations between amyloid burden and sensorimotor or motor-adjacent regions. The cohort was not uniformly CDR 0, and amyloid-positive participants were older and differed in sex distribution, so the structural result requires cautious interpretation. PET-defined resting-state motor-connectivity differences were smaller and did not survive false-discovery-rate correction, suggesting that resting state may be underpowered for motor-system effects.

\subsection*{Future directions}
Future studies should combine amyloid and tau biomarkers with age/sex-adjusted structural analyses, task-fMRI paradigms that directly engage manual dexterity, and quantitative finger-tapping or hand-function measurements.

\end{document}
